\section{Introduction}
\IEEEPARstart{T}{elevision} audiences no longer watch broadcasts in silence. During live debates, political talk shows, and reality competitions, viewers actively turn to second-screen platforms like Twitter to share commentary, memes, and screenshots \cite{SocialTV}. This stream of real-time viewer feedback offers broadcasters far richer signals than traditional audience ratings, providing immediate insight into public controversy, approval, and emotional engagement \cite{memorai}. 

Capturing these affective signals computationally, however, presents severe challenges. Social TV commentary is short, colloquial, and inherently multimodal. Viewers attach screengrabs and reaction images that do not merely mirror the written text; they exaggerate emotional tone, introduce sarcasm, or invert the literal message completely. 

The foundational question that motivated this investigation is: \textit{how much do attached images genuinely contribute to identifying human affective intent in multimodal social network posts?} During the manual annotation of our dataset, we observed a crucial empirical phenomenon: human annotator agreement increased substantially when annotators were provided with the attached images compared to reading the text alone. Visual context anchored ambiguous social commentary, establishing consensus where text alone yielded divergent interpretations. This observation directly motivates our central research hypothesis: \textit{does this visual grounding similarly carry over to neural models, or do visual features introduce semantic noise when discourse turns non-literal?}

Furthermore, while multilingual initiatives have advanced text-based emotion analysis, computational resources for Italian remain scarce, and existing benchmarks focus exclusively on unimodal text \cite{bianchi2021feel, emit, MultiEmotion}. To bridge this gap, we introduce \textbf{MultiEMiT}, a multimodal dataset of Italian Twitter commentary gathered from television broadcasts. MultiEMiT provides annotations across five pragmatic layers: fine-grained emotions under Plutchik's taxonomy \cite{Plutchik}, irony and sarcasm, offensiveness, target orientation (program container versus broadcast content), and semiotic text-image relationships \cite{MultiM}. Our benchmarking focuses on three primary classification tasks: Emotion Detection, Irony Detection, and Semiotic Relation Detection, with offensiveness and target serving as contextual signals. To expand training coverage for baseline text encoders, we also draw on the external unimodal EMiT (for Emotions) \cite{emit} and IronITA (for Irony) \cite{ironita2018} corpora.

An in-depth statistical analysis of MultiEMiT's multi-layered annotations revealed another compelling pattern: extensive empirical co-occurrences connect different annotation levels. For instance, irony correlates with joy, sarcasm pairs strongly with disgust and offensiveness, and program-directed commentary associates with viewer trust. Rather than treating these dimensions as isolated tasks or leaving neural models blind to these structural dependencies, we investigated whether artificial models could actively benefit from these cross-layer co-occurrences. To this end, we propose \textbf{Correlation Prompting}, a strategy that converts empirical Pearson correlation statistics directly into natural language prompt prefixes, providing explicit inductive priors during training and inference.

To evaluate these hypotheses systematically, we frame our study around three core research questions:
\begin{itemize}
    \item \textbf{RQ1 (Visual Context):} Does adding visual context help neural models detect emotion and irony, or do images introduce distracting noise in the classification of posts?
    \item \textbf{RQ2 (Correlation Prompting):} Does injecting statistical label correlations directly into text inputs improve accuracy across discriminative and generative architectures?
    \item \textbf{RQ3 (Architectural Paradigms):} How do compact, task-specialized discriminative models (BERT-ViT Dual-Encoders with Classifier Chains) compare with a 12 billion parameter foundation model (Pixtral 12B)?
\end{itemize}

The main contributions of this work are summarized as follows:
\begin{itemize}
    \item We introduce \textbf{MultiEMiT}, the first multi-layered multimodal dataset of Italian Social TV commentary, pairing text with broadcast media and reaction memes across five pragmatic dimensions.
    \item We empirically demonstrate the divergent impact of visual grounding across literal and figurative discourse: while images resolve textual ambiguity in emotion recognition, unguided visual features introduce semantic dissonance in irony detection.
    \item We propose \textbf{Correlation Prompting}, demonstrating that injecting multi-level statistical dependencies into model prompts effectively bridges cross-modal conflict and yields substantial performance gains.
    \item We present a rigorous comparative benchmark evaluating compact Dual-Encoders against the 12B parameter Pixtral model, identifying key failure modes such as exposure bias and hallucination loops in generative affective classification.
\end{itemize}

The rest of this paper is organized as follows. Section~\ref{sec:related_works} situates MultiEMiT within Italian affective computing, Social TV analytics, and multimodal document intent. Section~\ref{sec:dataset_analysis} describes the corpus collection, annotation protocol, and cross-dimensional correlations. Section~\ref{sec:experimental_design} defines our classification tasks, Dual-Encoder architecture, correlation prompting strategy, and Pixtral configurations. Section~\ref{sec:results} presents comparative evaluations and subset analyses across semiotic relations. Section~\ref{sec:conclusion} summarizes our findings and directions for future work.
